\section{Introduction}
Language-model research systems increasingly retrieve sources, synthesize claims, attach citations, and run automated verification \citep{aar2026,provenanceguard2026}. This creates a governance problem: a claim may be correct but attributed to the wrong source; a source may support a nearby claim but not the assigned one; a checker may be non-empty yet non-discriminating or share lineage with the answer; and an evaluator may be contaminated or tampered with after the fact. In each case, a scalar confidence score can remain high even though a prerequisite for scientific authority is absent.

We study a stricter object: a \emph{scientific-authority transition}. A proposal can be plausible and well supported but is promoted only when a registered set of content, provenance, checker, and evaluator prerequisites all pass. Unresolved prerequisites do not contribute a low score; they prevent promotion and leave the decision undetermined. This design is motivated by source-aware verification, multi-source attribution, claim-level auditability, evidence influence, iterative fact checking, evaluator-tampering benchmarks, and search-time contamination work \citep{provenanceguard2026,attributionbench2024,aar2026,provenai2026,fire2024,rewardhacking2026,stc2026,deepsciverify2026}. Recent academic-integrity and agentic-abstention benchmarks also make explicit that declining to complete an unsupported action can itself be the correct behavior rather than a failure to optimize task success \citep{sciintegrity2026,agentabstain2026}.

The contribution is deliberately narrower than those component techniques. We ask whether composing these requirements as a \emph{non-compensatory, protected authority transition} materially reduces false scientific-authority promotion while preserving legitimate clean coverage. Throughout, the pipeline built to that rule is called the \emph{governed pipeline}, so that each claim reads as a claim about the rule rather than about a product; Section~\ref{sec:methods} names the implementation that instantiates it.

The publication-authorizing campaign was executed only after the subject, comparator and ablation mechanisms, evaluator, statistical rules, and split-generation code were frozen. A separate 39-case live-model arm remains exploratory because post-run audit found label/denominator inconsistencies and a hidden-family execution leak; it is retained as adverse evidence and is not used to authorize, tune, or rescue the protected result.

Our main findings are: (i) the governed pipeline made 0/360 false promotions while the strongest frozen mechanism proxy made 180/360, with the paired interval entirely beyond the predeclared five-point margin; (ii) both promoted 60/60 clean positives, satisfying the clean-coverage non-inferiority safeguard; (iii) the registered undetermined-outcome contrast was not supported on that battery because every system was correct on a structurally saturated family, whereas a shortcut-resistant battery supported the same contrast at $1.0$ (95\% CI $[1.0,1.0]$) against the strongest comparator's 0/30 and at $0.5$ against the 15/30 of the one comparator mechanism that can escalate; this contrast measures the ability to express an undetermined outcome rather than finer-grained judgement; and (iv) every registered ablation increased false promotion without reducing clean coverage, with soft confidence producing the largest degradation.
